\begin{proposition}\label{prop:octic46}
No elliptic curve over an octic number field has a point of order $46$.
\end{proposition}
\begin{proof}
The rational Jacobian of $X_1(46)$ has rank zero by
\cite[Theorem~3.1(2)]{DEHMZ}, and its rational gonality is at least
$10$ by the finite-field bound in \cite[Table~1]{NV}.
Entirely cuspidal reduction at $3$ is excluded by
Lemma~\ref{lem:octic-odd-cuspidal}.

The complete noncuspidal closed-place counts over $\mathbb F_3$
through degree $8$ are
\[
 (0,0,0,55,44,55,264,781).
\]
They agree with an independent Frobenius similarity-class count.
The possible noncuspidal support types are one place of degree $4$
through $8$, or two quartic places, with repetition allowed.
All finite cusps of degree at most $4$ lift to rational cusps, so
every prime $q\in\{5,7,11,13,17,19\}$ is permitted for all
cuspidal complements.

The completed single-place Hecke branches give respectively
$5,4,5,24,71$ tests in degrees $4,5,6,7,8$, each with a
nonprincipal $A_q$ image for a permitted $q$.
For two quartic places, the first support runs through all five
diamond representatives and the second through all $55$ quartic
places. This includes all unordered pairs and repetition.

To accelerate these $275$ tests we also use
$\pi:X_1(46)\longrightarrow X_1(23)$, induced by $(E,P)\mapsto(E,2P)$.
By \cite[Corollary~4.14 and Table~4]{DEHMZ}, the rational Jacobian of
$X_1(23)$ is generated by cusps and has order $408991$.
At $3$ the exact degree-zero class group is $C_{9406793}$;
the image of the global rational cuspidal classes has order
$408991$, with quotient $C_{23}$. There is no cuspidal complement
for support type $(4,4)$, so Lemma~\ref{lem:octic-quotient-cusps}
requires its pushed-forward class to vanish in this quotient.
The quotient test, followed when needed by the permitted $A_q$
tests on $X_1(46)$, leaves no survivor among all $275$ pairs.
The point pushforward includes residue-field multiplicities.

These six support types exhaust every possible nonempty
noncuspidal reduction of a degree-$8$ divisor. The exact logs are
\texttt{N46\_p3\_single\_completed.log} and
\texttt{N46\_p3\_mixed\_quotient.log}; their coverage is checked
in \texttt{N46\_combined\_audit.json}. Together with the
entirely cuspidal branch this proves the proposition.
\end{proof}
